%------------------------------------------------------------------------------%
% Luis A. D'Afonseca
%
%------------------------------------------------------------------------------%
\documentclass[a4paper,12pt,fleqn]{article}

\usepackage{style_avaliacao}

%\ShowAnswers

\newcommand{\D}[2]{\frac{\partial {#1}}{\partial {#2}}}
\newcommand{\DS}[2]{\frac{\partial^2 {#1}}{\partial {#2}^2}}
\newcommand{\DM}[3]{\frac{\partial^2 {#1}}{\partial {#2} \partial {#3}}}

\newcommand{\vetor}[2]{\left(\begin{array}{c} #1 \\ #2 \end{array}\right)}

\newcommand{\veci}{\bm{i}}
\newcommand{\vecj}{\bm{j}}
\newcommand{\veck}{\bm{k}}

%------------------------------------------------------------------------------%
\begin{document}

\makeHeader{CFVVI}{Prova Substitutiva -- Turma 1}{11/09/2024}

\textbf{O exercício correspondente a prova que será substituida vale 50 pontos}

%P1
%Determine e esboce o domínio da função
%Encontre e esboce duas curvas de nível da função
%Calcule os limites solicitados, ou prove que o limite não existe
%
%P2
%Derivada direcional
%Regra da cadeia
%Linearização
%Plano tangente
%Pontos Críticos e regra derivada segunda
%Derivada função implicita
%
%P3
%Extremos em região fechada e limitada
%Multiplicadores de Lagrange
%Números Complexos

% 01
%------------------------------------------------------------------------------%
\exercise{25}
Calcule o limite
\(
  \lim_{(x, y) \to (0, 0)} \frac{x^2 - y^2}{x^2 + y^2}
\)
ou mostre que ele não existe.

\begin{answer}
  Restringindo a função à reta \( y = 0 \) temos
  \[
    \left(\lim_{(x, y) \to (0, 0)} \frac{x^2 - y^2}{x^2 + y^2}\right) \evalat{y=0}{}
    = \lim_{x\to 0} \frac{x^2}{x^2}
    = \lim_{x\to 0} 1
    = 1
  \]
  Porém, se restringirmos $f$ à reta \( x = 0 \), temos
  \[
    \left(\lim_{(x, y) \to (0, 0)} \frac{x^2 - y^2}{x^2 + y^2}\right) \evalat{x=0}{}
    = \lim_{y\to 0} \frac{-y^2}{y^2}
    = \lim_{y\to 0} -1
    = -1
  \]
  Os limites ao longo de diferentes caminhos são diferentes, logo o limite não existe.
  \clearpage
\end{answer}

% 02
%------------------------------------------------------------------------------%
\exercise{25}
Encontre e classifique os pontos críticos da função
\(
  f(x, y) = x^2 + y^2 - 4xy
\)

\begin{answer}
  Calculando as derivadas parciais de $f$
  \[
    f_x(x, y) = 2x - 4y, \qquad\qquad f_y(x, y) = 2y - 4x
  \]

  Como as derivadas parciais são polinômios elas existem em todos os pontos
  do plano. Portanto, temos apenas os pontos críticos onde as derivadas se anulam.
  Assim precisamos resolver o sistema \(\nabla f = 0\)
  \[
    \begin{cases}
      f_x(x, y) = 2x - 4y = 0 \\
      f_y(x, y) = 2y - 4x = 0
    \end{cases}
  \]
  que pode ser reescrito como
  \[
    \begin{cases}
      x = 2y \\
      y = 2x
    \end{cases}
  \]
  Substituindo \( y = 2x \) em \( x = 2y \), obtemos
  \[
    x = 4x
  \]
  portanto
  \[
    x = 0 \qquad \text{e} \qquad y = 0
  \]
  Encontramos o ponto crítico \((0, 0)\)

  Para classificar o ponto crítico, usamos o teste da segunda derivada.

  Calculando as derivadas de segundas de $f$
  \[
    f_{xx}(x, y) = 2 \qquad\qquad
    f_{yy}(x, y) = 2 \qquad\qquad
    f_{xy}(x, y) = -4
  \]

  O discriminante é
  \[
    D(x, y) = f_{xx}(x, y)f_{yy}(x, y) - \big(f_{xy}(x, y)\big)^2
    = 2 \times 2 - (-4)^2
    = 4 - 16
    = -12
  \]
  Como \(D(0, 0) = -12 <0\) o ponto \((0, 0)\) é um ponto de sela
  \clearpage
\end{answer}

% 03
%------------------------------------------------------------------------------%
\exercise{25}
Encontre os valores máximo e mínimo de \(f(x, y, z) = x - 2y + 5z\)
na esfera \(x^2+y^2+z^2=30\)

\begin{answer}
  Queremos os extremos da função
  \[
    f(x, y, z) = x - 2y + 5z
  \]
  restrita a
  \[
    g(x, y, z) = x^2 + y^2 + z^2 = 30
  \]
  Para aplicarmos multiplicadores de Lagrange precisamos das derivadas parciais
  das funções $f$ e $g$
  \[
    f_x(x, y, z) =  1 \qquad\qquad
    f_y(x, y, z) = -2 \qquad\qquad
    f_z(x, y, z) =  5
  \]
  \[
    g_x(x, y, z) = 2x \qquad\qquad
    g_y(x, y, z) = 2y \qquad\qquad
    g_z(x, y, z) = 2z
  \]
  Precisamos resolver o sistema \(\nabla f = \lambda \nabla g\) e \(g=30\)
  \[
    \begin{cases}
       \hphantom{-}1 = 2\lambda x \\
      -2 = 2\lambda y \\
       \hphantom{-}5 = 2\lambda z \\
       x^2 + y^2 + z^2 = 30
    \end{cases}
  \]
  Da primeiro equação temos
  \[
    \lambda = \frac{1}{2x}
  \]
  Isolando $y$ na segunda e substituindo $\lambda$ temos
  \[
    y
    = \frac{-2}{2\lambda}
    = \frac{-1}{\lambda}
    = -2x
  \]
  Fazendo o mesmo na terceira temos
  \[
    z
    = \frac{5}{2\lambda}
    = \frac{5\times 2x}{2}
    = 5x
  \]
  Substituindo $y$ e $z$ na quarta equação
  \begin{align*}
    x^2 + y^2 + z^2        & = 30 \\
    x^2 + (-2x)^2 + (5x)^2 & = 30 \\
    x^2 + 4x^2 + 25x^2     & = 30 \\
    30 x^2 = 30            &  \\
    x = \pm 1              &
  \end{align*}

  Para $x=1$
  \[
    y = -2x = -2 \qquad\qquad
    z = 5x = 5
  \]
  portanto
  \[
    f(1, -2, 5)
    = 1 - 2(-2) + 5\times 5
    = 1 + 4 + 25
    = 30
  \]

  Para $x=-1$
  \[
    y = -2x = 2 \qquad\qquad
    z = 5x = -5
  \]
  portanto
  \[
  f(-1, 2, -5)
    = -1 - 2\times 2 + 5(-5)
    = -1 - 4 - 25
    = -30
  \]

  O valor mínimo de $f$ é $-30$ e o valor máximo é $30$
\end{answer}

%------------------------------------------------------------------------------%
\end{document}
%------------------------------------------------------------------------------%
